\subsection{Definition of regular local rings}

Let A be a Noetherian local ring. Let \((x_{i})_{i\in I}\) be a family of elements of \(m_{A}\) . According to cor. 2 of prop. 4 of II, § 3, n° 2, it is equivalent to suppose that the family \((x_{i})_{i\in I}\) generates the ideal \(m_{A}\) of A, or that the classes of the \(x_{i}\) modulo \(m_{A}^{2}\) generate the \(\kappa_{A}\) -vector space \(m_{A}/m_{A}^{2}\) ; if this is so, one has dim(A) ≤ Card(I) by the scholium of § 3, n° 2. One therefore has the inequality

\[
\dim (A) \leqslant [ m _ {A} / m _ {A} ^ {2}: \kappa_ {A} ] \leqslant \operatorname{Card} (I)\tag{1}
\]

for every family \((x_{i})_{i\in I}\) generating the ideal \(m_{A}\) of A.

DEFINITION 1. --- The Noetherian local ring A is said to be regular if one has \(\dim(A) = [m_A/m_A^2:\kappa_A]\). One then calls a system of coordinates of A every family of elements of \(m_A\) whose classes modulo \(m_A^2\) form a basis of the \(\kappa_A\)-vector space \(m_A/m_A^2\).

A system of coordinates in a regular Noetherian local ring A is therefore a finite family \((x_{i})_{i\in I}\) generating the ideal \(m_{A}\) of A and such that \(\text{Card}(I) = \dim(A)\) . Conversely, if the ideal \(m_{A}\) of a Noetherian local ring A is generated by d elements with \(d \leqslant \dim(A)\) , the ring A is regular.

Examples. --- 1) The regular Noetherian local rings of dimension 0 (resp. 1) are the fields (resp. the discrete valuation rings) (VI, § 3, n° 6, prop. 9 and cor. 1 of th. 1 below). Let A be a discrete valuation ring; then an element t of m \(_{A}\) is a uniformizer if and only if \{t\} is a system of coordinates of A.

\begin{enumerate}
\def\labelenumi{\arabic{enumi})}
\setcounter{enumi}{1}
\tightlist
\item
  Let \(k\) be a field and let \(n\) be a positive integer. The ring of formal power series \(k[[X_1, ..., X_n]]\) is a regular Noetherian local ring of dimension \(n\) (§ 3, n° 4, cor. 3 of prop. 8). Let \(F_1, ..., F_n\) be formal power series without constant term in \(k[[X_1, ..., X_n]]\) ; for the sequence \((F_1, ..., F_n)\) to be a system of coordinates of \(k[[X_1, ..., X_n]]\) , it is necessary and sufficient that the matrix \(\left(\frac{\partial F_i}{\partial X_j}(0, ..., 0)\right)\) be invertible.
\end{enumerate}

More generally, let A be a regular Noetherian local ring of dimension \(r\) ; then \(A[[X_1, ..., X_n]]\) is a regular Noetherian local ring of dimension \(r + n\) . If \((a_1, ..., a_r)\) is a system of coordinates of A, then \((a_1, ..., a_r, X_1, ..., X_n)\) is a system of coordinates of \(A[[X_1, ..., X_n]]\) .

\begin{enumerate}
\def\labelenumi{\arabic{enumi})}
\setcounter{enumi}{2}
\tightlist
\item
  Let A be a complete regular Noetherian local ring of dimension \(r\) . The ring A \(\{X_1, ..., X_n\}\) of restricted formal power series is a regular Noetherian local ring of dimension \(r + n\) (§ 3, n° 4, remark 1). If \((a_1, ..., a_r)\) is a system of coordinates of A, then \((a_1, ..., a_r, X_1, ..., X_n)\) is a system of coordinates of A \(\{X_1, ..., X_n\}\) .
\end{enumerate}

* 4) Let \(k\) be a complete non-discrete valued field. The ring of formal power series in \(n\) variables which converge in a neighborhood of 0 in \(k^n\) is a regular Noetherian local ring of dimension \(n\) ( \(\S\) 3, n° 4, remark 2). *

\begin{enumerate}
\def\labelenumi{\arabic{enumi})}
\setcounter{enumi}{4}
\tightlist
\item
  Let \(k\) be a field, A a \(k\) -algebra of finite type which is an integral domain, and \(\mathfrak{m}\) a maximal ideal of A. The Noetherian local ring \(\mathrm{A}_{\mathfrak{m}}\) is regular if and only if one has \(\dim(\mathrm{A}) = [\mathfrak{m}/\mathfrak{m}^{2}:\mathrm{A}/\mathfrak{m}]\) : indeed, one has \(\dim(\mathrm{A}_{\mathfrak{m}}) = \dim(\mathrm{A})\) ( \(\S 2\) , \(no. 4\) , cor. 2 of th. 3) and the vector spaces \(\mathfrak{m}/\mathfrak{m}^{2}\) and \(\mathrm{mA}_{\mathfrak{m}}/(\mathrm{mA}_{\mathfrak{m}})^{2}\) over the field A/m are isomorphic (II, \(\S 3\) , \(no. 3\) , prop. 9). In particular, if \(k\) is algebraically closed, the stated condition is equivalent to \(\dim(\mathrm{A}) = [\mathfrak{m}/\mathfrak{m}^{2}:k]\) (V, \(\S 3\) , \(no. 3\) , prop. 1).
\end{enumerate}

* 6) Let X be an algebraic variety over a perfect field \(k\) . Then X is nonsingular at a point \(x\) if and only if the local ring of X at \(x\) is regular. *

* 7) Let A be a regular Noetherian local ring. It will be seen later that the Noetherian local ring \(A_{p}\) is regular for every prime ideal p of A. *

PROPOSITION 1.--- Let A and B be Noetherian local rings and \(\rho:A\to B\) a local homomorphism making B a flat A-module. Suppose that one has \(m_{B} = B \cdot \rho(m_{A})\) . Then one has \(\dim(A) = \dim(B)\) and B is regular if and only if A is regular.

The first assertion follows from cor. 1 of prop. 7 of § 3, no. 4. Since B is flat over A, one may identify \(\mathfrak{m}_{\mathrm{B}}^{k} = \mathbf{B}.\rho (\mathfrak{m}_{\mathrm{A}}^{k})\) with \(\mathbf{B}\otimes_{\mathbf{A}}\mathfrak{m}_{\mathbf{A}}^{k}\) for every integer \(k\geqslant 0\) , hence \(\mathfrak{m}_{\mathrm{B}} / \mathfrak{m}_{\mathrm{B}}^{2}\) with \(\mathbf{B}\otimes_{\mathbf{A}}(\mathfrak{m}_{\mathrm{A}} / \mathfrak{m}_{\mathrm{A}}^{2})\) or again with \(\kappa_{\mathrm{B}}\otimes_{\kappa_{\mathrm{A}}}(\mathfrak{m}_{\mathrm{A}} / \mathfrak{m}_{\mathrm{A}}^{2})\) . One therefore has

\[
\left[ \mathfrak {m} _ {\mathrm{B}} / \mathfrak {m} _ {\mathrm{B}} ^ {2}: \kappa_ {\mathrm{B}} \right] = \left[ \mathfrak {m} _ {\mathrm{A}} / \mathfrak {m} _ {\mathrm{A}} ^ {2}: \kappa_ {\mathrm{A}} \right],\tag{2}
\]

whence the proposition immediately.

COROLLARY. --- A Noetherian local ring A is regular if and only if its completion \(\hat{A}\) is.

Indeed, \(\hat{\mathbf{A}}\) is flat over \(\mathbf{A}\) , and one has \(\mathfrak{m}_{\hat{\mathbf{A}}} = \hat{\mathbf{A}}. \mathfrak{m}_{\mathbf{A}}\) (III, § 3, no. 4, th. 3 and § 2, no. 12, cor. 2 to prop. 16).

